% Source: https://artofproblemsolving.com/community/c3h1849707_a_right_triangle
% Stable identity: aops-topic-1849707-q01
Use coordinates with $C=(0,0)$, $A=(b,0)$, $B=(0,a)$, where $a=BC$, $b=CA$ and $c=AB=\sqrt{a^2+b^2}$. The line $AB$ is $ax+by=ab$, and the vector $(a,b)$ is perpendicular to $\overrightarrow{AB}=(-b,a)$, has length $c$, and points away from $C$. So the square outside the triangle has
\[
Q=A+(a,b)=(a+b,\,b),\qquad P=B+(a,b)=(a,\,a+b).
\]

\begin{asy}
usepackage("amsmath");
size(5.5cm);
defaultpen(fontsize(10pt));
real a=3,bb=4;
pair C=(0,0),A=(bb,0),B=(0,a),Q=A+(a,bb),P=B+(a,bb);
pair Y=extension(A,B,C,Q),X=extension(A,B,C,P);
draw(A--B--C--cycle,linewidth(1));
draw(A--Q--P--B,linewidth(1));
draw(C--Q,blue);draw(C--P,red);
draw((0.35,0)--(0.35,0.35)--(0,0.35),linewidth(0.5));
dot(X);dot(Y);
label("$C$",C,SW);label("$A$",A,S);label("$B$",B,W);label("$Q$",Q,E);label("$P$",P,N);
label("$X$",X,dir(195),red);label("$Y$",Y,dir(268),blue);
label("\mbox{}",(min(currentpicture,user=true)+max(currentpicture,user=true))/2,black); // keeps the text after the figure black when LaTeX typesets the labels inline
\end{asy}

Points of $AB$ can be written $A+s(B-A)=(b-sb,\,sa)$, where $s$ is the fraction of the way from $A$ to $B$, so that the distance from $A$ is $sc$. Let $D=a^2+ab+b^2$.
\begin{itemize}
\item $Y$ is on $CQ$: $Y=t(a+b,b)$ with $a\,t(a+b)+b\cdot tb=ab$, so $t=\dfrac{ab}{D}$. Its second coordinate $\dfrac{ab^2}{D}$ equals $sa$, so $s_Y=\dfrac{b^2}{D}$.
\item $X$ is on $CP$: $X=u(a,a+b)$ with $a\cdot ua+b\,u(a+b)=ab$, so $u=\dfrac{ab}{D}$. Its second coordinate $\dfrac{ab(a+b)}{D}$ equals $sa$, so $s_X=\dfrac{b(a+b)}{D}$.
\end{itemize}
Since $s_Y<s_X$, the order along $AB$ is $A,Y,X,B$, and
\[
AY=\dfrac{b^2c}{D},\qquad XY=(s_X-s_Y)c=\dfrac{abc}{D},\qquad XB=(1-s_X)c=\dfrac{(D-ab-b^2)c}{D}=\dfrac{a^2c}{D}.
\]
Therefore
\[
XY^2=\dfrac{a^2b^2c^2}{D^2}=\dfrac{b^2c}{D}\cdot\dfrac{a^2c}{D}=AY\cdot XB .\qquad\blacksquare
\]
